\section{Allowing Intelligence to resume command}

\paragraph{Introducing Albert Gleizes}

\begin{quotex}
Intellectualism, the abuse of reason, has destroyed man's kingdom. The cause of the present trouble is rather a sickness of the mind than a breakdown of the body: the organs, abandoned to their own devices, soon become atrophied or hypertrophied and rapidly bring to naught the organism's last attempts at resistance; \emph{no wonder, therefore, that a few intellectually-gifted people are beginning to wake up---and none too soon}.
\end{quotex}

Some thoughts from the French artist Albert Gleizes\footnote{\url{https://en.wikipedia.org/wiki/Albert\_Gleizes}}, who brought the thoughts of Rene Guenon into the world of Cubism.

\begin{quotationx}
We dream of progress and our dream is becoming a nightmare. A dream implies a loss of lucidity brought on by great fatigue and the numbing of consciousness. We had good cause for feeling fatigued. For centuries we worked in the manner of men; even while in decline we struggled to preserve our nobility. Exhausted, we began to lose our heads; our waking hours were haunted by dreams; only the illusion of progress kept us alive. It was leading us peacefully towards death, stimulating our failing strength, supplying us with orthopaedic material. We cannot complain if is has served its purpose in bringing us to our goal.

If among the intellectuals there are some who accept death and do not seek salvation, they are at liberty to do so, but those who wish to be reborn must purify their souls.

To the despairing impulse which sees in the conduct of individuals or groups nothing more than subconscious or hyperconscious reflexes, is it not possible forthwith to oppose the value of conscience, the forcefulness of common sense, the happiness which springs from measured and direct action such as permits man to control himself and to claim the fruits of the tasks he has shouldered? The very fact that a contrast is possible, implies that both attitudes are legitimate; let us not deny to the former its rights, but rather let us accept the duties of the latter. In any case, they differ only in the different combination and application of identical values which, according to their point of departure, may either be heralds of death or bringers of life. Impenitent, intellectualism converts the dreams into a deadly nightmare; repentant, it awakes from its dreaming and allows Intelligence to resume command.
\end{quotationx}



\flrightit{Posted on 2011-01-17 by Cologero}

\begin{center}* * *\end{center}

\begin{footnotesize}
\begin{sffamily}
\texttt{Graham on 2011-01-17 at 18:22 said:}

Real intellectuality, real depth, is simple, as poetry is simple. Practice can mean doing and feeling in a measured, attentive mood. Such a practice is infectious -- from moment to moment and from person to person -- and renovates `everyday' life.

In order to allow this a spell of solitude and danger may be necessary. More than one, perhaps.

\hfill

\texttt{Max on 2015-02-28 at 11:10 said:}

``Cubism, regarded as a teaching, is actually the way that reliably leads to the simple truth, that is with maximal precision obtained by minimal means. The strictly personal vision expressed as condensed as possible.''

Moving in the same circles as Gleizes was Ivan Aguéli who also worked together with Guénon and is worth looking up for someone interested in the relation of traditional and modern art. He published some thoughts on ``the pure art'' in French in the magazine La Gnose (Pages dédiées à Mercure), also available in Swedish translation here:

\url{http://www.agueliportalen.se/2010/05/17/den-rena-konsten/}

For those who read French or Italian, here is a page that contains scanned originals mainly from La Gnose and Il Convito, for example an article titled ``El Malamatiyah'' on a practice recently mentioned on Gornahoor:
\url{http://agueliportalen.startaenweb.com/wp-content/blogs.dir/35/files/}


``And since `the portrait of a space' -if i dare say so -only lends itself to be expressed by exact proportion between distance, lines and luminosity, it follows that cubism above all is the uncompromising architecture working with the light and dark integrities of empty space. We consequently see that what is unusual with cubism is just its daring. It is nothing but the strictest method -at least today -known for studying the exactness of the hues and the precision of drawing.''

Only the un-initiated would reject cubism for its novelty since it is in fact said to represent a purification. Cubism has to be understood as an effort to explore the relation of simultainity and succession or maybe eternity to the present moment. It should be mentioned that Aguéli always sought to unify complementary opposites, in his painting for example the harmonization of line and colour in their common principle to bring forth depth. He also claims that an Artist is one who God created a ``personal sun'' for and that he should strive to understand and express this light under the pretext that to the extent one is oneself, freedom is realized.

A translation from the original Swedish, written in 1912 by Aguéli:

``A healthy mind, if so only equipped with common sense and the most elementary concepts of humanity will find out that it is equally in her nature to be something positively exceptional and unique as to be anybody. And to conclude the view, as the social problems solution is situated in the various forms of reunification of the two contrary currents, that is the harmony between the individual and the common good, so does the solution to the questions of the soul lie in the individuals capacity to first live his own life, then others life, so everyones life, and finally, everythings life. But it is of utmost importance to never forget that the first and basic condition to reach the height of this transcendental pantheism is to first and foremost live ones own life; since during this long and curious evolution, one does not for one moment stop being oneself. On the contrary, since the intensity of life is made into the sole purpose of life.

(... )

And so does we come to the conclusion that the respectively individual and common developments, evolutions and revolutions are fundamentally inseperable and subjugating eachother mutually. To be free one has to be among free men, since no one is free amongst slaves, not even the slave owner. As one can not achieve one without the other it goes without saying that they must be achieved simultaneously. To wish one after the other is already done and completed seems to me an absurdity.''


\hfill

\end{sffamily}
\end{footnotesize}

